\section{Integral Extensions}

\begin{definition}[Integral element]
    Let \(S\) be an integral domain, and let \(R\subseteq S\) be a subring containing
    \(1_S\). Let \(\alpha\in S\). We say that \(\alpha\) is \emph{integral over \(R\)}
    if there exists a monic polynomial
    \[
        f(X)\in R[X]
    \]
    such that
    \[
        f(\alpha)=0.
    \]
    Equivalently, there exist \(a_0,a_1,\dots,a_{n-1}\in R\) such that
    \[
        \alpha^n+a_{n-1}\alpha^{n-1}+\cdots+a_1\alpha+a_0=0.
    \]
\end{definition}

\begin{proposition}[Equivalent characterizations of integrality]
    Let \(S\) be an integral domain, and let \(R\subseteq S\) be a subring containing
    \(1_S\). For \(\alpha\in S\), the following conditions are equivalent:
    \begin{enumerate}
        \item \(\alpha\) is integral over \(R\).
        \item \(R[\alpha]\) is a finitely generated \(R\)-module.
        \item There exists a subring \(T\) of \(S\) such that
              \[
                  R\subseteq T,\qquad \alpha\in T,
              \]
              and \(T\) is finitely generated as an \(R\)-module.
    \end{enumerate}
\end{proposition}

\begin{proof}
    We prove the implications separately.

    \medskip

    \noindent
    \((2)\Rightarrow (3)\).
    Take
    \[
        T=R[\alpha].
    \]
    Then \(T\) is a subring of \(S\), contains \(R\), contains \(\alpha\), and is finitely
    generated as an \(R\)-module by assumption.

    \medskip

    \noindent
    \((1)\Rightarrow (2)\).
    Assume that \(\alpha\) is integral over \(R\). Then there exists a monic polynomial
    \[
        f(X)=X^n+a_{n-1}X^{n-1}+\cdots+a_1X+a_0\in R[X]
    \]
    such that \(f(\alpha)=0\). Hence
    \[
        \alpha^n+a_{n-1}\alpha^{n-1}+\cdots+a_1\alpha+a_0=0,
    \]
    so
    \[
        \alpha^n
        =
        -a_{n-1}\alpha^{n-1}
        -\cdots
        -a_1\alpha
        -a_0.
    \]
    Thus \(\alpha^n\) is an \(R\)-linear combination of
    \[
        1,\alpha,\alpha^2,\dots,\alpha^{n-1}.
    \]
    Multiplying this relation by powers of \(\alpha\), we see inductively that every
    power \(\alpha^m\) with \(m\ge n\) is also an \(R\)-linear combination of
    \[
        1,\alpha,\alpha^2,\dots,\alpha^{n-1}.
    \]
    Therefore every element of \(R[\alpha]\) is an \(R\)-linear combination of
    \[
        1,\alpha,\alpha^2,\dots,\alpha^{n-1}.
    \]
    Hence
    \[
        R[\alpha]
        =
        R\cdot 1+R\alpha+\cdots+R\alpha^{n-1},
    \]
    so \(R[\alpha]\) is a finitely generated \(R\)-module.

    \medskip

    \noindent
    \((3)\Rightarrow (1)\).
    Assume that there exists a subring \(T\) of \(S\) such that
    \[
        R\subseteq T,\qquad \alpha\in T,
    \]
    and \(T\) is finitely generated as an \(R\)-module.

    Let
    \[
        \varepsilon_1,\dots,\varepsilon_n
    \]
    be a set of generators of \(T\) as an \(R\)-module. Since \(\alpha\in T\) and \(T\)
    is a subring, for each \(i\) we have
    \[
        \alpha \varepsilon_i\in T.
    \]
    Hence there exist elements \(a_{ji}\in R\) such that
    \[
        \alpha\varepsilon_i
        =
        \sum_{j=1}^n a_{ji}\varepsilon_j,
        \qquad 1\le i\le n.
    \]
    Let
    \[
        A=(a_{ji})\in M_n(R).
    \]
    Writing
    \[
        \varepsilon=
        \begin{pmatrix}
            \varepsilon_1 \\
            \vdots        \\
            \varepsilon_n
        \end{pmatrix},
    \]
    the above relations can be written as
    \[
        (\alpha I_n-A)\varepsilon=0.
    \]
    Multiplying both sides by the adjugate matrix
    \[
        \operatorname{adj}(\alpha I_n-A),
    \]
    we obtain
    \[
        \det(\alpha I_n-A)\varepsilon=0.
    \]
    Thus
    \[
        \det(\alpha I_n-A)\varepsilon_i=0
        \qquad
        \text{for all } i=1,\dots,n.
    \]
    Since \(\varepsilon_1,\dots,\varepsilon_n\) generate \(T\) as an \(R\)-module, it follows that
    \[
        \det(\alpha I_n-A)\,T=0.
    \]
    But \(1\in T\), so
    \[
        \det(\alpha I_n-A)=0.
    \]
    Now define
    \[
        p(X)=\det(XI_n-A)\in R[X].
    \]
    This is a monic polynomial of degree \(n\). Since
    \[
        p(\alpha)=\det(\alpha I_n-A)=0,
    \]
    we conclude that \(\alpha\) is integral over \(R\).
\end{proof}

\begin{lemma}[Finite generation of a generated subring]
    Let \(S\) be an integral domain, and let \(R\subseteq S\) be a subring containing
    \(1_S\). Let \(M_1,M_2\) be subrings of \(S\) such that
    \[
        R\subseteq M_i\subseteq S,
        \qquad i=1,2.
    \]
    If \(M_1\) and \(M_2\) are finitely generated as \(R\)-modules, then the subring
    \[
        R[M_1,M_2]\subseteq S
    \]
    generated by \(R,M_1,M_2\) is finitely generated as an \(R\)-module.
\end{lemma}

\begin{proof}
    Suppose
    \[
        M_1=R\alpha_1+\cdots+R\alpha_m
    \]
    as an \(R\)-module, and
    \[
        M_2=R\beta_1+\cdots+R\beta_n
    \]
    as an \(R\)-module.

    We claim that \(R[M_1,M_2]\) is generated as an \(R\)-module by the finite set
    \[
        \{\alpha_i\beta_j\mid 1\le i\le m,\ 1\le j\le n\}.
    \]
    Set
    \[
        N=\sum_{i=1}^m\sum_{j=1}^n R\alpha_i\beta_j.
    \]
    Clearly \(N\subseteq R[M_1,M_2]\). We show that \(N\) is a subring of \(S\) containing
    \(M_1\) and \(M_2\).

    Since \(1\in M_2\), every \(\alpha_i\in M_1\) can be written as
    \[
        \alpha_i=\alpha_i\cdot 1\in N.
    \]
    Similarly, since \(1\in M_1\), every \(\beta_j\in M_2\) belongs to \(N\). Hence
    \[
        M_1\subseteq N,
        \qquad
        M_2\subseteq N.
    \]

    It remains to show that \(N\) is closed under multiplication. Take two generators
    \(\alpha_i\beta_j\) and \(\alpha_{i'}\beta_{j'}\). Then
    \[
        (\alpha_i\beta_j)(\alpha_{i'}\beta_{j'})
        =
        (\alpha_i\alpha_{i'})(\beta_j\beta_{j'}).
    \]
    Since \(M_1\) and \(M_2\) are subrings, we have
    \[
        \alpha_i\alpha_{i'}\in M_1,
        \qquad
        \beta_j\beta_{j'}\in M_2.
    \]
    Therefore there exist \(r_k,s_\ell\in R\) such that
    \[
        \alpha_i\alpha_{i'}
        =
        \sum_{k=1}^m r_k\alpha_k,
        \qquad
        \beta_j\beta_{j'}
        =
        \sum_{\ell=1}^n s_\ell\beta_\ell.
    \]
    Thus
    \[
        (\alpha_i\alpha_{i'})(\beta_j\beta_{j'})
        =
        \sum_{k=1}^m\sum_{\ell=1}^n r_ks_\ell\,\alpha_k\beta_\ell
        \in N.
    \]
    Hence \(N\) is closed under multiplication.

    Therefore \(N\) is a subring of \(S\) containing \(M_1\) and \(M_2\). By minimality of
    \(R[M_1,M_2]\), we get
    \[
        R[M_1,M_2]\subseteq N.
    \]
    Since the reverse inclusion is clear, we have
    \[
        R[M_1,M_2]=N.
    \]
    Hence \(R[M_1,M_2]\) is finitely generated as an \(R\)-module.
\end{proof}

\begin{theorem}[Integral elements form a subring]
    Let \(S\) be an integral domain, and let \(R\subseteq S\) be a subring containing
    \(1_S\). Then the set of integral elements of \(S\) over \(R\) forms a subring of
    \(S\) containing \(R\).
\end{theorem}

\begin{proof}
    Let \(\alpha,\beta\in S\) be integral over \(R\). By the characterization of
    integral elements, \(R[\alpha]\) and \(R[\beta]\) are finitely generated
    \(R\)-modules. Hence, by the finite-generation lemma for generated subrings,
    \[
        R[\alpha,\beta]
    \]
    is also a finitely generated \(R\)-module.

    Since
    \[
        \alpha+\beta,\quad \alpha-\beta,\quad \alpha\beta
        \in R[\alpha,\beta],
    \]
    again by the characterization of integrality, the elements
    \[
        \alpha+\beta,\qquad \alpha-\beta,\qquad \alpha\beta
    \]
    are integral over \(R\).

    Moreover, every element \(r\in R\) is integral over \(R\), since it satisfies the
    monic polynomial
    \[
        X-r\in R[X].
    \]
    Therefore the integral elements of \(S\) over \(R\) form a subring of \(S\)
    containing \(R\).
\end{proof}

\begin{definition}[Algebraic integer]
    Let \(\mathbb Z\subseteq \mathbb C\). If \(z\in\mathbb C\) is integral over
    \(\mathbb Z\), then \(z\) is called an \emph{algebraic integer}.
\end{definition}

\begin{definition}[Integrally closed subring]
    Let \(S\) be an integral domain, and let \(R\subseteq S\) be a subring containing
    \(1_S\). We say that \(R\) is \emph{integrally closed in \(S\)} if every element
    of \(S\) which is integral over \(R\) already lies in \(R\).
\end{definition}

\begin{definition}[Integrally closed domain]
    Let \(R\) be an integral domain with quotient field \(K=\operatorname{Frac}(R)\).
    We say that \(R\) is \emph{integrally closed} if \(R\) is integrally closed in
    its quotient field \(K\). Equivalently, every element of \(K\) integral over
    \(R\) belongs to \(R\).
\end{definition}

\begin{theorem}
    Every unique factorization domain is integrally closed.
\end{theorem}

\begin{proof}
    Let \(R\) be a unique factorization domain, and let
    \[
        K=\operatorname{Frac}(R)
    \]
    be its quotient field. We show that \(R\) is integrally closed in \(K\).

    Let \(\alpha\in K\) be integral over \(R\). Write
    \[
        \alpha=\frac{b}{a},
    \]
    where \(a,b\in R\), \(a\neq 0\), and \(\gcd(a,b)=1\).

    Since \(\alpha\) is integral over \(R\), there exist \(n\geq 1\) and
    \(c_0,c_1,\dots,c_{n-1}\in R\) such that
    \[
        \alpha^n+c_{n-1}\alpha^{n-1}+\cdots+c_1\alpha+c_0=0.
    \]
    Substituting \(\alpha=b/a\), we get
    \[
        \left(\frac{b}{a}\right)^n
        +c_{n-1}\left(\frac{b}{a}\right)^{n-1}
        +\cdots
        +c_1\left(\frac{b}{a}\right)
        +c_0
        =0.
    \]
    Multiplying by \(a^n\), we obtain
    \[
        b^n+c_{n-1}b^{n-1}a+\cdots+c_1ba^{n-1}+c_0a^n=0.
    \]
    Thus
    \[
        b^n
        =
        -a\left(c_{n-1}b^{n-1}
        +c_{n-2}b^{n-2}a
        +\cdots
        +c_1ba^{n-2}
        +c_0a^{n-1}\right).
    \]
    Hence
    \[
        a\mid b^n.
    \]

    If \(a\) is not a unit, then \(a\) has a prime divisor \(p\). Since \(a\mid b^n\),
    we have
    \[
        p\mid b^n.
    \]
    Because \(p\) is prime in the UFD \(R\), it follows that
    \[
        p\mid b.
    \]
    But \(p\mid a\) as well, contradicting the assumption that \(\gcd(a,b)=1\).

    Therefore \(a\) must be a unit in \(R\). Hence
    \[
        \alpha=\frac{b}{a}\in R.
    \]
    So every element of \(K\) integral over \(R\) belongs to \(R\). Therefore \(R\)
    is integrally closed.
\end{proof}

